\chapter{Giới thiệu}
\label{chap:introduction}

\section{Bối cảnh và động lực}

Trong những năm gần đây, số lượng thiết bị IoT được triển khai trong các
nhà máy, tòa nhà thông minh và hệ thống giám sát môi trường tăng theo
cấp số nhân. Mỗi thiết bị liên tục phát ra các bản ghi cảm biến --- nhiệt
độ, độ ẩm, độ rung, áp suất --- với tần suất mili-giây đến vài giây.
Theo các khảo sát gần đây về IoT Big Data~\cite{iot_big_data_survey_2022},
một hệ thống cảm biến quy mô công nghiệp có thể dễ dàng sinh ra hàng
chục triệu đến hàng tỷ bản ghi mỗi ngày, đồng thời đòi hỏi độ trễ xử
lý đầu cuối chỉ tính bằng giây để phục vụ cảnh báo và ra quyết định.

Đặc thù này khiến các mô hình xử lý dữ liệu truyền thống --- lưu trữ
tập trung, xử lý batch theo lô --- gặp hai nghịch lý:

\begin{itemize}
  \item Kích thước dữ liệu lớn khiến việc lưu trữ toàn bộ rồi xử lý
        theo lô trở nên tốn kém và chậm;
  \item Giá trị của dữ liệu suy giảm nhanh theo thời gian, đòi hỏi xử
        lý liên tục gần với thời điểm phát sinh.
\end{itemize}

Hệ quả là các hệ thống IoT ở quy mô lớn phải được xây dựng theo mô hình
\textbf{streaming}: dữ liệu được đưa vào pipeline ngay khi sinh ra,
được xử lý bằng các engine phân tán, và được lưu trữ ở dạng có cấu trúc
cho phép truy vấn cả real-time lẫn lịch sử~\cite{stream_processing_evolution_2024}.

Ở Việt Nam, một số nghiên cứu gần đây cũng bắt đầu áp dụng các công
nghệ Big Data cho bài toán giám sát giao thông và chất lượng không khí
tại TP.HCM, cho thấy nhu cầu thực tiễn của mô hình này trong bối cảnh
smart-city. Báo cáo này xây dựng trên cùng nguyên lý nhưng chọn bài
toán cảm biến IoT tổng quát làm use-case để có tính đại diện cao.

\section{Mục tiêu nghiên cứu}

Khác với một bài triển khai đơn thuần ("cài Kafka + Spark + Delta rồi
chạy thử"), báo cáo lấy các câu hỏi về \textit{hệ thống} làm trọng tâm.
Năm câu hỏi nghiên cứu (Research Question, viết tắt RQ) được đặt ra:

\begin{description}
  \item[RQ-A (Throughput \& Latency):] Throughput và ingestion lateness
        (trễ từ khi hiện tượng xảy ra đến khi nền tảng nhận được event)
        của pipeline Kafka--Spark thay đổi như thế nào khi input rate
        tăng từ 1k đến 50k event/giây?
  \item[RQ-B (Scalability):] Kafka partition count và số Spark executor
        ảnh hưởng thế nào đến khả năng mở rộng theo chiều ngang?
  \item[RQ-C (Fault Recovery):] Checkpointing trên HDFS và Delta transaction
        log giúp hệ thống phục hồi thế nào sau khi một worker hoặc
        executor bị kill?
  \item[RQ-D (Data Veracity):] Các cơ chế veracity --- kiểm tra required-field
        và physical bounds, deduplication hai pass, watermark và metric
        lateness --- xử lý các IoT event đến muộn, trùng lặp và hỏng hiệu quả
        đến đâu; record accounting (chấp nhận / cách ly / loại công khai) có
        được bảo toàn không; và một chính sách thích nghi theo lateness
        (đề xuất ở mức khái niệm) có thể mang lại lợi ích gì so với cấu hình
        cố định?
  \item[RQ-E (Storage \& Query):] Phân vùng Delta và OPTIMIZE/-- VACUUM
        ảnh hưởng thế nào đến storage efficiency và analytical query latency?
\end{description}

Năm RQ này biến đề tài từ "dùng Spark để tính nhiệt độ trung bình" thành
một đánh giá mang bản chất Big Data Systems, nơi các đặc trưng Volume,
Velocity, Veracity và tính phân tán được đo lường một cách có hệ thống.

\section{Phạm vi}

Báo cáo giới hạn phạm vi như sau:

\begin{itemize}
  \item \textbf{Stack công nghệ:} Apache Kafka 4.3, Apache Spark 4.2 với
        Structured Streaming, Delta Lake, HDFS 3.5, Apache Airflow,
        Prometheus, Grafana, JMX Exporter.
  \item \textbf{Môi trường triển khai:} Docker Compose trên một máy phát
        triển, gồm 1 Kafka broker, 1 Spark master + 2 worker, 1 NameNode
        + 2 DataNode, và các dịch vụ hỗ trợ (Airflow, Prometheus, Grafana).
        Khi benchmark fault tolerance, có thể chuyển sang 3 Kafka broker
        nếu tài nguyên cho phép.
  \item \textbf{Dữ liệu:} IoT sensor event với schema đơn giản gồm 9 trường
        \texttt{event\_id}, \texttt{sensor\_id}, \texttt{event\_time},
        \texttt{ingest\_time}, \texttt{sensor\_type}, \texttt{value},
        \texttt{unit}, \texttt{location}, \texttt{sequence\_no}
        (định nghĩa hình thức tại Mục~\ref{subsec:event-tuple}). Nguồn dữ liệu
        có thể là (i) sensor simulator tự sinh với tải có kiểm soát, hoặc
        (ii) replay từ dataset Intel Berkeley~\cite{intel_berkeley_dataset}.
  \item \textbf{Ngoài phạm vi:} Kubernetes ở mức production (Strimzi,
        Spark on K8s, Airflow Helm chart), HDFS high-availability với
        QJM, TLS/SASL + ACL đầy đủ, edge/fog computing. Đây được xem
        là các hướng phát triển trong Chương~\ref{chap:conclusion}.
\end{itemize}

\section{Đóng góp của báo cáo}

Báo cáo đóng góp bốn điểm chính:

\begin{enumerate}
  \item \textbf{Kiến trúc tham chiếu} cho hệ thống Big Data streaming IoT
        với ngăn xếp hoàn toàn mã nguồn mở và có tài liệu chính thức
        đồng bộ;
  \item \textbf{Khung khái niệm veracity-aware event-time}: một mô hình hình
        thức gồm định nghĩa event và phân loại thời gian, lateness, ngữ nghĩa
        watermark, veracity vector, metric rolling, chính sách FAST/CAUTIOUS
        và vòng đời công bố aggregate, kèm các bất biến và giả thuyết kiểm
        chứng --- một mô hình tổng hợp có thể kiểm chứng cho pipeline
        Kafka--Spark--Delta--HDFS, không tuyên bố tính mới của từng khái niệm
        riêng lẻ;
  \item \textbf{Phương pháp benchmark} gồm các kịch bản (smoke, load ramp,
        executor/worker scaling, partition scaling, fault injection,
        veracity injection, storage replay, compaction) với các metric được
        định nghĩa rõ ràng và mapping tới từng RQ cùng từng giả thuyết;
  \item \textbf{Số liệu minh họa} dựa trên các benchmark công bố gần đây
        (\cite{shufflebench_2024, fault_recovery_stream_2024}) và tài
        liệu kỹ thuật của Confluent, được sử dụng để minh hoạ
        phương pháp phân tích thay vì đo lường trực tiếp trên cluster
        thật. Phần adaptive veracity được trình bày như đề xuất/giả thuyết
        chưa triển khai, chưa có số đo.
\end{enumerate}

\section{Cấu trúc báo cáo}

Báo cáo được tổ chức như sau:

\begin{itemize}
  \item Chương~\ref{chap:background} trình bày cơ sở lý thuyết và các
        công nghệ chính;
  \item Chương~\ref{chap:implementation} mô tả kiến trúc hệ thống, schema
        dữ liệu và cấu hình triển khai;
  \item Chương~\ref{chap:results} trình bày phương pháp benchmark và phân
        tích số liệu minh họa;
  \item Chương~\ref{chap:conclusion} tổng kết và đề xuất hướng phát triển.
\end{itemize}